\section{第~\ref{sec:motorlearning}~章　运动学习}

最小均方规则和单独误差导线的优势是定理；误差导线线路的结构、同时出现时的削弱、资格迹与延迟的匹配、后效以及技能的自发恢复都已测量；整条线路作为有监督学习工作并建立内部模型，是主流假说；双过程模型的数字是按本书模型计算的。

{\footnotesize
\setlength{\tabcolsep}{3pt}\renewcommand{\arraystretch}{1.15}
\begin{longtable}{>{\raggedright}p{0.5cm}>{\raggedright}p{4.0cm}>{\raggedright}p{5.6cm}>{\raggedright}p{2.3cm}>{\raggedright\arraybackslash}p{1.7cm}}
\toprule
序号 & 命题 & 实验及测量内容 & 文献 & 状态 \\
\midrule\endfirsthead
\toprule
序号 & 命题 & 实验及测量内容 & 文献 & 状态 \\
\midrule\endhead
1 & 规则~\eqref{eq:lms}平均而言沿误差平方的梯度下降 & 自适应滤波器理论的结果：与输入和输出误差成正比的修正，就是对均方误差梯度的随机下降。 & \cite{WidrowHoff1960} & 理论 \\
2 & 每个输出有误差导线时，速度与输出数目无关；只有一个公共数时，速度随噪声神经元数目下降（命题~\ref{prop:teachers}） & 线性网络的学习曲线：按神经元的随机扰动学习，比精确梯度慢的倍数等于被扰动神经元的个数。 & \cite{Werfel2005} & 理论 \\
3 & 误差导线线路作为有监督学习工作（命题~\ref{prop:errorwire}） & 由线路结构推出的理论。第7--10行的实验与之相符；整条线路正是这样工作，尚未证明。 & \cite{Marr1969, Albus1971} & 假说 \\
4 & 扩展层：约$7\cdot10^{10}$个小型\nt{GLUT}神经元，比大脑其余部分加起来还多 & 在溶解组织的悬液中计数细胞核：人小脑有$69\cdot10^9$个神经元，全脑为$86\cdot10^9$个。体视学：老年男性小脑中有$101\cdot10^9$个小神经元。 & \cite{Azevedo2009, Andersen1992cereb} & 已测量 \\
5 & 升高输入维数使所需的关系变为线性 & 关于划分数目的定理：$n$个输入的带阈值加权和能实现约$2n$个一般位置向量的任意标注。小脑正是利用了这一点，属于第3行假说的一部分。 & \cite{Cover1965} & 理论 \\
6 & 约$3\cdot10^7$个输出\nt{GABA}神经元，每个约有$10^5$个输入 & 人小脑体视学：$30.5\cdot10^6$个输出神经元。电子显微镜，大鼠：每个输出神经元约有$1.75\cdot10^5$个来自扩展层的输入。输出神经元的抑制性递质见综述。 & \cite{Andersen1992cereb, NapperHarvey1988, Ito2001} & 已测量 \\
7 & 每个输出神经元恰好一根误差导线，约每秒一次，每次引起强烈爆发 & 记录综述：每个输出神经元接受一个攀缘\nt{GLUT}输入，它以约$1$\,Hz的频率引起复合脉冲。 & \cite{Ito2001} & 已测量 \\
8 & 爆发的概率取决于动作误差的方向 & 猴，小脑动眼区，扫视适应：视觉误差的方向决定复合脉冲的概率，误差大小决定其时间；爆发改变下一次尝试中的简单脉冲，并使动作沿该细胞偏好的误差方向偏移。 & \cite{Herzfeld2018} & 已测量 \\
9 & 误差爆发与活跃输入同时出现会削弱该输入（$-\eta e_i x_j$） & 实验综述：同时刺激扩展层输入和误差导线，引起该输入的长时程削弱；单独刺激一个输入则不会。 & \cite{Ito2001} & 已测量 \\
10 & 资格迹的长度与误差延迟匹配：延迟$100\ms$时，被削弱最多的是$100\ms$前活跃的输入 & 切片和活体小鼠：在负责动眼学习的区域，削弱被精确地调节到输入与误差之间约$120\ms$的间隔——正是这项任务中误差信号的传送时间；在另一区域，不同细胞调节到不同的间隔。 & \cite{Suvrathan2016} & 已测量 \\
11 & 线路是学习身体内部模型的自适应滤波器~\eqref{eq:adaptivefilter} & 由线路结构推出的模型。尚无把学到的滤波器作为身体传递函数来测量的直接实验。 & \cite{Fujita1982} & 假说 \\
12 & 预测减去自身动作的后果，所以不能给自己挠痒 & 人：自己的触摸比别人同样的触摸感觉更不痒；在扫描仪中，躯体感觉皮层对自己的触摸响应较弱，而当动作触及皮肤时小脑的活动较小。用减去预测来解释是假说。 & \cite{Blakemore1998} & 假说 \\
13 & 在外力作用下偏差减小，撤去外力后手臂偏向另一侧 & 人在力场中握着机器人手柄伸手：轨迹恢复为直线；撤去力场后，后效几乎是最初畸变的镜像；它还会迁移到未经训练的空间区域。 & \cite{ShadmehrMussaIvaldi1994} & 已测量 \\
14 & 两个过程在学习~\eqref{eq:tworate}；反向扰动之后，技能会自行恢复 & 人，力场，然后是短暂的反向力场和误差被钳制的试次：修正自行回到原来的值；双过程模型能再现这一点，单过程模型不能。 & \cite{Smith2006} & 已测量 \\
15 & 图~\ref{fig:tworate}的数字：$200$次动作后为$0.84$（慢过程$0.63$）；$6$次反向；快过程$-0.53$，慢过程$0.46$；$40$次动作内回升到$0.37$ & 按\texttt{models/\allowbreak motor\_\allowbreak learning.py}计算，$A_f=0.92$，$B_f=0.1$，$A_s=0.996$，$B_s=0.02$：$0.840$（$0.634$和$0.205$），$6$次动作，$-0.531$和$0.457$，第$40$次动作时峰值$0.371$。 & \texttt{models/\allowbreak motor\_\allowbreak learning.py} & 模型 \\
16 & 需要两个过程，是因为身体以不同速度变化 & 理论：在具有多种寿命的干扰下，最优估计给出多个时间常数不同的滤波器，并解释了自发恢复和加速的再学习。 & \cite{Kording2007} & 假说 \\
\bottomrule
\end{longtable}}
